%------------------------------------------------------------------------------%
\documentclass[fleqn,utf8,aspectratio=169,12pt]{beamer}

\usepackage{calculo_varias_variaveis-1}

\title{Derivadas Parciais e Continuidade}

%------------------------------------------------------------------------------%
\begin{document}

\addtitle

%------------------------------------------------------------------------------%
\section{Derivadas Parciais e Continuidade}

%------------------------------------------------------------------------------%
\begin{frame}{Derivada e Continuidade em Uma Variável}
  No cálculo de funções de \emph{uma variável}

  \pause
  \medskip
  se uma função é \emph{derivável} em um ponto $a$

  \pause
  \medskip
  ela é \emph{contínua} nesse ponto
\end{frame}

%------------------------------------------------------------------------------%
\begin{frame}{Derivadas Parciais e Continuidade}
  Se as derivadas parciais de $f$ existem em $(a, b)$

  \pause
  \emph{não podemos dizer nada} sobre a continuidade de $f$ no ponto
\end{frame}

%------------------------------------------------------------------------------%
\begin{question}
  Considerando a função
  \(\;
    f(x, y) = \begin{cases}
      0, & xy \neq 0 \\
      1, & xy = 0
    \end{cases}
  \) \\[5mm]
  \begin{enumerate}
    \item Mostre que suas derivadas parciais existem na origem \\[3mm]
    \item Prove que $f$ não é contínua na origem
  \end{enumerate}
\end{question}

%------------------------------------------------------------------------------%
\begin{resolution}{Mostrar que as Derivadas Parciais Existem}
  Fazendo $y=0$,
  \pause
  temos $xy=0$,
  \pause
  portanto \(f(x, 0) = 1\)
  \pause
  e
  \[
    \D{f}{x}\evalat{y=0}{}
    =\; \D{1}{x}
    \;=\; 0
  \]

  \pause
  \bigskip
  Fazendo $x=0$,
  \pause
  temos $xy=0$,
  \pause
  portanto \(f(0, y) = 1\)
  \pause
  e
  \[
    \D{f}{y}\evalat{x=0}{}
    =\; \D{1}{y}
    \;=\; 0
  \]

\end{resolution}

%------------------------------------------------------------------------------%
\begin{resolution}{Mostrar que $f$ é Descontínua na Orígem}
  Verificando a condição de existência do limite na origem

  \pause
  Considerando a reta $y=0$, com $x\neq 0$
  \pause
  \[
    f(x, 0)
    \pause
    = 1
  \]

  \pause
  Considerando a reta $y=x$, com $x\neq 0$
  \[
    \pause
    f(x, x)
    \pause
    = 0
  \]

  \pause
  Como $f$ tende a valores diferentes por
  caminhos diferentes para a origem,\\
  o limite não existe
\end{resolution}

%------------------------------------------------------------------------------%
\framedgraphic{Derivada e Continuidade}{Derivadas_parciais_continuidade}{}

%------------------------------------------------------------------------------%


%------------------------------------------------------------------------------%
\section{Lista Mínima}

%------------------------------------------------------------------------------%
\begin{frame}{Lista Mínima}
  Cálculo Vol. 2 do Thomas 12\textsuperscript{a} ed. -- Seção 14.3
  \begin{enumerate}
    \item Estudar o texto da seção
  \end{enumerate}

  \vfill
  Atenção: A prova é baseada no livro, não nas apresentações
\end{frame}

%------------------------------------------------------------------------------%
\end{document}
%------------------------------------------------------------------------------%
